\chapter{Ergebnisse}
\label{cha:ergebnisse}

Dieses Kapitel berichtet die erhobenen Daten zunächst rein darstellend: den deskriptiven Überblick
über Autorenschaft, Reflexionsstufen und Survey-Kennwerte (Abschnitt~\ref{sec:ergebnisse}), die
Illustration des Kategoriensystems an realen Entscheidungssituationen
(Abschnitt~\ref{sec:illustration-kategoriensystem}) sowie die Gegenüberstellung von Selbstauskunft und
dokumentiertem Bearbeitungsverhalten (Abschnitt~\ref{sec:selbstauskunft-verhalten}). Die
Interpretation dieser Befunde und die Beantwortung der Forschungsfragen erfolgen anschließend in
Kapitel~\ref{cha: diskussion}.

\section{Deskriptiver Überblick}
\label{sec:ergebnisse}

Dieser Abschnitt stellt die erhobenen Daten zunächst rein deskriptiv dar. Ihre inhaltliche
Interpretation und die Beantwortung der Forschungsfragen erfolgen in Kapitel~\ref{cha: diskussion}.
Angesichts der kleinen Stichprobe sind alle Werte deskriptiv zu lesen. Auf Inferenzstatistik wird
bewusst verzichtet (vgl. Abschnitt~\ref{sec:datenanalyse}). Die vollständigen Tabellen finden sich in
\hyperref[app:kodierungstabelle]{Anhang~A.2} bis \hyperref[app:statistik]{A.6}.

\subsection*{Autorenschaft der Bearbeitungen}
Von den 31 kodierten Analyseeinheiten ist die Autorenschaft nur bei 11 eindeutig studentisch.
Tabelle~\ref{tab:ergebnisse-autorenschaft} schlüsselt die Verteilung nach Gruppe und nach Sicherheit
der Einstufung auf. Fasst man die Kategorien \glqq unklar\grqq{} und \glqq KI-generiert\grqq{}
zusammen, liegt der Anteil der Bearbeitungen ohne eindeutig studentische Autorenschaft bei Team AI bei
80\,\% und beim Reflexionsbot bei 50\,\%. Betrachtet man dagegen nur die gesichert KI-generierten
Fälle, fällt der Unterschied zwischen beiden Gruppen deutlich geringer aus.

\begin{table}[htb]
\centering
\caption{Autorenschaft der kodierten Bearbeitungen nach Gruppe}
\label{tab:ergebnisse-autorenschaft}
\renewcommand{\arraystretch}{1.3}
\setlength{\tabcolsep}{8pt}
\begin{tabular}{>{\raggedright\arraybackslash}p{5.2cm}
                >{\raggedleft\arraybackslash}p{3.4cm}
                >{\raggedleft\arraybackslash}p{3.6cm}}
\toprule
\textbf{Autorenschaft} & \textbf{Team AI (n=15)} & \textbf{Reflexionsbot (n=16)} \\
\midrule
studentisch  & 3 & 8 \\
unklar       & 4 & 1 \\
KI-generiert & 8 & 7 \\
\midrule
nicht eindeutig studentisch & 12 (80\,\%) & 8 (50\,\%) \\
\bottomrule
\end{tabular}
\end{table}

\subsection*{Reflexionsstufen der studentisch verfassten Bearbeitungen}
Für die Fälle mit eindeutig studentischer Autorenschaft verteilen sich die Reflexionsstufen wie in
Abbildung~\ref{fig:ergebnisse-stufen} dargestellt. Deskriptiv liegt der Mittelwert der Stufen bei Team
AI mit 1.33 (n=3) über dem des Reflexionsbots mit 0.75 (n=8); die Mediane liegen bei 2.0 gegenüber 0.5.
Bei drei gegenüber acht auswertbaren Fällen ist dieser Unterschied jedoch nicht belastbar. Bereits ein
einzelner abweichender Fall verschiebt den Mittelwert der kleineren Gruppe um mehr als 0,3 Stufen.

Kein einziger Fall der gesamten Stichprobe erreicht mit bestätigter studentischer Autorenschaft
Stufe 4. Texte, die zunächst wie kritische Reflexion wirkten, stellten sich beim Abgleich mit den
Rohdaten durchgehend als KI-generiert oder unklar heraus
(vgl. Abschnitt~\ref{sec:illustration-kategoriensystem}). Kritische Reflexion im Sinne von Hatton und
Smith ist allerdings auch unabhängig vom eingesetzten Werkzeug ein seltenes Phänomen, insbesondere bei
Studierenden im ersten Semester und in einer einmaligen Übungseinheit. Das vollständige Ausbleiben
dieser Stufe ist daher für sich genommen kein Befund über die verglichenen KI-Systeme.

\begin{figure}[htb]
\centering
\begin{tikzpicture}
\begin{axis}[
	ybar,
	bar width=16pt,
	width=0.85\textwidth,
	height=5.4cm,
	symbolic x coords={Stufe 0, Stufe 1, Stufe 2, Stufe 3, Stufe 4},
	xtick=data,
	xlabel={Reflexionsstufe nach Hatton und Smith},
	ylabel={Anzahl Fälle},
	xlabel style={font=\small},
	ylabel style={font=\small},
	x tick label style={font=\small},
	y tick label style={font=\small},
	ymin=0, ymax=5,
	enlarge x limits=0.15,
	legend style={at={(0.5,-0.30)}, anchor=north, legend columns=2, font=\small, draw=none,
	              /tikz/every even column/.append style={column sep=0.6cm}},
	nodes near coords,
	nodes near coords style={font=\scriptsize},
	axis lines*=left,
	ytick={0,1,2,3,4,5},
]
\addplot[fill=codeblue!70, draw=codeblue] coordinates {(Stufe 0,4) (Stufe 1,3) (Stufe 2,0) (Stufe 3,1) (Stufe 4,0)};
\addplot[fill=codegreen!60, draw=codegreen] coordinates {(Stufe 0,1) (Stufe 1,0) (Stufe 2,2) (Stufe 3,0) (Stufe 4,0)};
\legend{Reflexionsbot (n=8), Team AI (n=3)}
\end{axis}
\end{tikzpicture}
\caption{Verteilung der Reflexionsstufen nach Gruppe (nur Fälle mit eindeutig studentischer
Autorenschaft)}
\label{fig:ergebnisse-stufen}
\end{figure}

\subsection*{Objektiver Wissenstest}
Im Wissenstest zu Scrum und Prozessmodellen (0 bis 6 Punkte) zeigt der Vergleich nach Systemgruppen
keinen Vorteil des Reflexionsbots. Die Reflexionsbot-Gruppe fällt von Prä M=5.56 auf Post M=4.44
(n=9), während die Generic-AI-Gruppe mit Prä und Post M=5.43 (n=7) stabil bleibt
(Tab.~\ref{tab:survey-wissenstest}). Etwa die Hälfte dieses Rückgangs entfällt auf einen einzelnen
Fall, der im Post-Survey alle sechs Fragen falsch beantwortete. Ohne diesen Fall liegt die
Reflexionsbot-Gruppe bei Prä M=5.62 und Post M=5.00 (n=8); ein leichter Rückgang bleibt damit auch
ohne ihn bestehen.

Die Aussagekraft dieses Vergleichs ist allerdings durch einen Deckeneffekt begrenzt. Beide Gruppen
liegen bereits im Prä-Test nahe der Höchstpunktzahl von rund 5,5 von 6 Punkten, sodass sich
Lernzuwächse kaum abbilden lassen. Aus dem Test ist damit weder ein Systemvorteil noch dessen Fehlen
belastbar abzuleiten. Dieser Vorbehalt gilt für jeden Vergleich, der auf diesem Instrument aufbaut,
also auch für die im Folgenden berichtete Aufteilung nach dem Bearbeitungsverhalten.

Teilt man dieselben Personen nach dem tatsächlichen Bearbeitungsverhalten auf, erreichen eigenständig
arbeitende Personen im Post-Test M=5.75 (n=4) gegenüber M=4.50 (n=10) bei Personen mit ganz oder
teilweise ausgelagerter Bearbeitung (vgl. Abschnitt~\ref{sec:selbstauskunft-verhalten}). Dieser
Abstand ist jedoch nicht als Lerneffekt zu lesen. Alle vier eigenständig arbeitenden Personen
erreichten bereits im Prä-Test die volle Punktzahl (Prä M=6.00, SD=0.00), sodass bei ihnen rechnerisch
kein Zuwachs mehr möglich war; ihr Post-Wert entspricht einem leichten Rückgang von $-$0.25. Die
Vergleichsgruppe startete bereits im Prä-Test unterhalb der Höchstpunktzahl und fiel ebenfalls ab. Der
Unterschied im Post-Test spiegelt damit überwiegend eine schon vor der Übung bestehende
Leistungsdifferenz wider und nicht einen unterschiedlichen Lernzuwachs.

\subsection*{Selbsteingeschätztes Prozessverständnis und Systembewertung}
Bei den Items zum reinen Beschreiben sinkt die Selbsteinschätzung in beiden Gruppen leicht. Bei den
Items zum Begründen und Abwägen bleibt die Reflexionsbot-Gruppe stabil oder verbessert sich leicht,
mit einer Veränderung von Post minus Prä zwischen 0.00 und $+$0.33. In der Generic-AI-Gruppe sinkt
sie dagegen durchgehend, mit Werten zwischen $-$0.43 und $-$0.86
(Tab.~\ref{tab:survey-prozessverstaendnis}). Beim selbstberichteten Reflexionsverhalten sinkt die
Selbsteinschätzung nach der Übung in beiden Gruppen
(Tab.~\ref{tab:survey-reflexionsverhalten}).

Vier Post-Survey-Items bilden die Grundlage für die Beantwortung von TF3 in
Abschnitt~\ref{sec:beantwortung-forschungsfragen} und sind in
Abbildung~\ref{fig:survey-vergleich} gegenübergestellt. Beim Item zur Kompetenzentwicklung
(\glqq Die Arbeit an den Dilemma-Situationen hat meine Fähigkeit verbessert, Prozessentscheidungen zu
begründen\grqq{}) liegt die Reflexionsbot-Gruppe mit M=3.10 leicht über der Generic-AI-Gruppe mit
M=2.80. Deutlicher fällt der Unterschied beim Item \glqq war hilfreicher als erwartet\grqq{} aus
(M=3.40 gegenüber M=2.20). Umgekehrt ist die Bereitschaft zur erneuten Nutzung beim Reflexionsbot
niedriger (M=2.50 gegenüber M=3.20), und Reflexionsbot-Nutzende hätten häufiger das jeweils andere
System bevorzugt (M=3.30 gegenüber M=2.20).

\begin{figure}[htb]
\centering
\begin{tikzpicture}
\begin{axis}[
	ybar,
	bar width=14pt,
	width=0.68\textwidth,
	height=5.7cm,
	symbolic x coords={Kompetenzentwicklung, Hilfreicher als erwartet, Erneute Nutzung, Anderes System bevorzugt},
	xtick=data,
	% Ungedrehte, zweizeilig gesetzte Beschriftungen: gedreht ueberlappten sie
	% sich bei dieser Achsenbreite.
	xticklabels={{Kompetenz-\\entwicklung}, {Hilfreicher\\als erwartet},
	             {Erneute\\Nutzung}, {Anderes System\\bevorzugt}},
	xticklabel style={font=\small, align=center},
	ylabel={Zustimmung (1--5)},
	ylabel style={font=\small},
	y tick label style={font=\small},
	ymin=0,
	ymax=5,
	enlarge x limits=0.15,
	% Legende rechts neben dem Diagramm statt darunter: spart Bauhoehe und
	% haelt die Beschriftung vom gedrehten x-Achsen-Text getrennt.
	legend style={at={(1.03,1)}, anchor=north west, legend columns=1, font=\small,
	              draw=none, cells={anchor=west}},
	nodes near coords,
	nodes near coords style={font=\scriptsize},
	axis lines*=left,
	ytick={0,1,2,3,4,5},
]
\addplot[fill=codeblue!70, draw=codeblue] coordinates {
	(Kompetenzentwicklung,3.10) (Hilfreicher als erwartet,3.40) (Erneute Nutzung,2.50) (Anderes System bevorzugt,3.30)
};
\addplot[fill=codegreen!60, draw=codegreen] coordinates {
	(Kompetenzentwicklung,2.80) (Hilfreicher als erwartet,2.20) (Erneute Nutzung,3.20) (Anderes System bevorzugt,2.20)
};
\legend{Reflexionsbot, Generisches System}
\end{axis}
\end{tikzpicture}
\caption[Vergleich zentraler Survey-Mittelwerte]{Vergleich zentraler Survey-Mittelwerte zwischen
Reflexionsbot- und Generic-AI-Gruppe (jeweils n=10)}
\label{fig:survey-vergleich}
\end{figure}

\subsection*{Kognitive Last}
Die Skala zur kognitiven Last (1 = very easy bis 5 = very difficult) ergibt kein einheitliches Bild
(Tab.~\ref{tab:survey-kognitive-last}). Bei den drei Teilaufgaben und bei den beiden Items zum Umgang
mit dem KI-System liegt die Reflexionsbot-Gruppe zwischen 0.00 und 0.40 Skalenpunkten höher, am
deutlichsten beim Bewerten der Gegenargumente des Systems (M=3.10 gegenüber M=2.80). Die
zusammenfassende Einschätzung des gesamten kognitiven Aufwands fällt dagegen \emph{gegenläufig} aus.
Hier liegt die Generic-AI-Gruppe mit M=3.20 über der Reflexionsbot-Gruppe mit M=3.00. Angesichts von
Standardabweichungen zwischen 1.00 und 1.66 bei zehn Personen je Gruppe sind sämtliche dieser
Unterschiede kleiner als die Streuung innerhalb der Gruppen. Aus dieser Skala lässt sich daher weder
ableiten, dass der Reflexionsbot als anstrengender empfunden wurde, noch das Gegenteil.

\section{Illustration des Kategoriensystems an realen Entscheidungssituationen}
\label{sec:illustration-kategoriensystem}

Um zu verdeutlichen, wie die Reflexionsstufen nach Hatton und Smith
\cite{hattonSmithReflectionTeacherEducation1995} auf das vorliegende Material angewendet wurden, wird
im Folgenden je ein reales Beispiel pro Stufe herangezogen. Alle Beispiele stammen von Personen, deren
Autorenschaft nach Abgleich mit den Rohdaten (Chatlogs, Zusatzdokumente, Formulierungsvergleiche
zwischen Personen) als eindeutig studentisch eingestuft wurde. Sie stammen bewusst nicht alle aus
derselben Entscheidungssituation, da die genauere Prüfung der Rohdaten mehrere ursprünglich als
studentisch angenommene Texte als KI-generiert entlarvt hat (vgl. Abschnitt~\ref{sec:datenanalyse}).

\textbf{Stufe 0, keine eigene Reflexion:} Ein Teilnehmer im Team Reflexionsbot antwortet auf wiederholte
Rückfragen nur mit: \glqq Stop asking questions!!!\grqq{} und liefert zu keinem Zeitpunkt eine eigene
inhaltliche Antwort.

\textbf{Stufe 1, deskriptives Schreiben (Descriptive Writing):} Eine Teilnehmerin im Team Reflexionsbot begründet ihre Priorisierung
auf mehrfache Rückfrage des Bots hin nur mit: \glqq Health > trustworthiness\grqq{}. Die Aussage bleibt
eine bloße Rangfolge ohne inhaltliche Auseinandersetzung mit dem Gegenargument des Systems.

\textbf{Stufe 2, deskriptive Reflexion (Descriptive Reflection):} Eine Teilnehmerin bei Team AI begründet ihre Wahl von
Datensatz A in Situation B des Szenarios SmartDose so: \glqq A gives better reliability but might run
into problems when it is later on revealed the data cannot be used and the ai has to be
retrained\grqq{}. Hier wird ein Kriterium (Zuverlässigkeit) benannt und einem konkreten Risiko
gegenübergestellt, jedoch ohne die nicht gewählte Option weiter zu erwägen.

\textbf{Stufe 3, dialogische Reflexion (Dialogic Reflection):} Ein Teilnehmer im Team Reflexionsbot wägt in Situation C des
Szenarios LearnLoop einen Nachteil explizit gegen einen Vorteil ab: \glqq While there is a risk of
context switching, this approach minimizes idle time and maintains momentum.\grqq{} Der Nachteil wird
benannt und bewusst gegen den Nutzen der eigenen Entscheidung abgewogen.

\textbf{Stufe 4, kritische Reflexion (Critical Reflection):} Für diese Stufe konnte in der gesamten Stichprobe kein Fall mit
bestätigter studentischer Autorenschaft gefunden werden. Mehrere Texte erfüllten die formalen
Kriterien der Stufe augenscheinlich sehr gut, etwa eine vollständige Abwägung von Datenschutz,
Zulassungsfähigkeit und Modellqualität unter ausdrücklichem Verweis auf GDPR und MDR. Beim Abgleich
mit den zugehörigen Chatlogs und Zusatzdokumenten stellten sich jedoch ausnahmslos alle diese Texte
als KI-generiert heraus, unter anderem, weil zwei unterschiedliche Personen fast wortgleiche
Formulierungen einreichten oder weil ein Link zum zugrunde liegenden ChatGPT-Verlauf mitgeliefert
wurde. Diese Beobachtung wird in Abschnitt~\ref{sec:beantwortung-forschungsfragen} als eigenständiger
Befund aufgegriffen: sprachlich und strukturell überzeugende Antworten waren in dieser Stichprobe ein
stärkeres Indiz für KI-Autorenschaft als für tiefe eigene Reflexion.

Das Beispiel zeigt, dass Reflexionstiefe in den vorliegenden Daten nicht in erster Linie vom
verwendeten KI-System abhing, sondern davon, ob eine Person sich überhaupt eigenständig mit der
Entscheidung auseinandersetzte, statt die Aufgabe zu vermeiden oder an die KI zu delegieren. Dass diese Bearbeitungshaltung eher
personengebunden als szenarienabhängig ist, zeigt sich daran, dass die Autorenschaft bei 10 der 11
Personen mit zwei Bearbeitungen über beide Szenarien identisch ausfällt. Nur in einem Fall wechselt
sie von \glqq studentisch\grqq{} zu \glqq unklar\grqq{} (vgl. \hyperref[app:kodierungstabelle]{Anhang~A.2}). Die
vollständige Kodierung aller Fälle sowie die Häufigkeitsverteilung je Gruppe sind im Anhang
dokumentiert.

\section{Selbstauskunft und dokumentiertes Bearbeitungsverhalten}
\label{sec:selbstauskunft-verhalten}

Der Abgleich der Autorenschafts-Einstufung aus Abschnitt~\ref{sec:datenanalyse} mit den
Post-Survey-Antworten derselben Personen (verknüpft über den persönlichen Code, vgl.
Abschnitt~\ref{sec:fragebogen-experiment}) zeigt eine weitere Facette des zentralen Befunds dieser
Arbeit. Von 14 Personen, für die sowohl eine Autorenschafts-Einstufung als auch Survey-Daten vorliegen,
bearbeiteten vier ihre Szenarien durchgehend eigenständig. Bei den übrigen zehn war mindestens eine
Bearbeitung ganz oder teilweise an die KI ausgelagert.

Im Wissenstest zu Scrum und Prozessmodellen schnitt die eigenständig arbeitende Gruppe deutlich besser
ab (Mittelwert 5.75 von 6 Punkten) als die Gruppe mit ausgelagerter Bearbeitung (Mittelwert 4.50 von 6
Punkten). Bei der Selbsteinschätzung des eigenen Reflexionsverhaltens im Survey zeigte sich dagegen
kein vergleichbarer Unterschied zwischen beiden Gruppen.

Besonders deutlich wird das Missverhältnis bei der wahrgenommenen Nützlichkeit des jeweiligen
Systems: Eine Person, die in der qualitativen Kodierung selbst angab, die Aufgaben nur kopiert und das
System nie aktiv genutzt zu haben, bewertete es im Survey trotzdem als sehr nützlich für die eigene
Reflexion. Umgekehrt bewerteten mehrere Personen mit eindeutig eigenständiger, aber wenig reflektierter
Bearbeitung dasselbe System als wenig hilfreich. Selbstauskunft im Fragebogen allein hätte den
Unterschied zwischen echter Auseinandersetzung und vollständiger Auslagerung an die KI in diesen
Fällen nicht sichtbar gemacht. Erst der Abgleich mit den Rohdaten der Bearbeitung deckte ihn auf. Diese
kleine Stichprobe (n=4 gegenüber n=10) erlaubt keine statistische Verallgemeinerung, liefert aber ein
inhaltliches Argument dafür, Selbstauskunft und dokumentiertes Verhalten in zukünftigen Studien zu
diesem Thema grundsätzlich getrennt zu erheben und gegeneinander zu prüfen.
